Part I asked whether a small interactive program can be shipped as a neural network held to its behaviour
specification, and showed that it can, one interface at a time. This paper changes the question from
generalisation to exactness, reports the result, and plans the step from interfaces to machines. Trained on the
whole of its spec, one network per interface is exact: of the \nPlanned{} interfaces we set out to train,
\nNewDemos{} of them rebuilt from the demo set of declarative dragging (a sliding puzzle, Tower of Hanoi with a
solver given only by example, a kanban board, Sokoban, a meeting scheduler), every finite network that has
finished (\nExactSpecs{} of \nFiniteSpecs{} as this draft is built) reproduces every transition with no rule
broken, on every seed, at \minParamsK--\maxParamsK k parameters and minutes to an hour on a local GPU;
restricting a network's answers to the states the spec's invariants allow removes rule breaks without
retraining. The recipe stops at an
enumeration wall, and we propose to go around it rather than through it: machines wired from exact blocks. We
define blocks with ports, three couplings (independent, event-wired, read-wired) with their training cost,
argue, with a proof sketch, that exactness composes under correct wiring, and specify the block architecture (heads chosen by field
type, a proposer for large state spaces, a compiled core with trained holes), the runtime (a synchronous tick,
bindings for keyboard and mouse, selection and undo, drags inside a running machine, discrete hand-offs between
continuous blocks) and an evaluation protocol for machines. Five demos are registered as experiments with
predictions and falsifiers: an odometer of digit networks, a self-running chain checked at every hand-off, a
machine you can grab mid-run, a desktop card board under keyboard and mouse, and a ball rolling across track
blocks. Everything measured is generated from the scripts and logs committed with this paper.
